\documentclass[a4paper, serif, 10pt]{moderncv}

%% moderncv
% style and color
\moderncvstyle{banking}
\moderncvcolor{black}

% renew moderncv command to adapt for CJK colon
\renewcommand*{\cvitem}[3][.25em]{%
  \ifstrempty{#2}{}{\hintstyle{#2}：}{#3}%
  \par\addvspace{#1}}

\renewcommand*{\cvitemwithcomment}[4][.25em]{%
  \savebox{\cvitemwithcommentmainbox}{\ifstrempty{#2}{}{\hintstyle{#2}：}#3}%
  \setlength{\cvitemwithcommentmainlength}{\widthof{\usebox{\cvitemwithcommentmainbox}}}%
  \setlength{\cvitemwithcommentcommentlength}{\maincolumnwidth-\separatorcolumnwidth-\cvitemwithcommentmainlength}%
  \begin{minipage}[t]{\cvitemwithcommentmainlength}\usebox{\cvitemwithcommentmainbox}\end{minipage}%
  \hfill% fill of \separatorcolumnwidth
  \begin{minipage}[t]{\cvitemwithcommentcommentlength}\raggedleft\small\itshape#4\end{minipage}%
  \par\addvspace{#1}}

%% page layout/margins
\usepackage[top=2.5cm, bottom=2.5cm, left=1.5cm, right=1.5cm]{geometry}

\nopagenumbers{}

%% Babel config for English language
\usepackage[english]{babel}

%% fontspec
\usepackage{fontspec}

\IfFontExistsTF{Linux Libertine}{
  \setmainfont[Ligatures={TeX, Common}, Numbers=Lining, ItalicFont=Linux Libertine]{Linux Libertine}
}{}
\IfFontExistsTF{Linux Libertine O}{
  \setmainfont[Ligatures={TeX, Common}, Numbers=Lining, ItalicFont=Linux Libertine O]{Linux Libertine O}
}{}

%% CTeX
% CJK support, used to show CJK characters in the resume
%
% - fontset=none: disable builtin fontset but instead set the CJK font manually
% - heading=false: disable ctex heading
% - punct=kaiming: use kaiming punctuations styles for CJK
% - scheme=plain: use plain scheme, do not override \normalsize font size
% - space=auto: space settings for CJK characters
%
% ref:
% - http://ctan.mirrorcatalogs.com/language/chinese/ctex/ctex.pdf
\usepackage[UTF8, heading=false, punct=kaiming, scheme=plain, space=auto]{ctex}

\IfFontExistsTF{Noto Serif CJK SC}{
  \setCJKmainfont{Noto Serif CJK SC}
}{}
\IfFontExistsTF{Noto Sans CJK SC}{
  \setCJKsansfont{Noto Sans CJK SC}
}{}

%% line spacing
\usepackage{setspace}
\setstretch{1.125}

%% URL styling - use normal text instead of monospace
\urlstyle{same}

% Auto-underline all links
% Defer until \begin{document} so hyperref is already loaded
\AtBeginDocument{%
  \let\oldhref\href
  \renewcommand{\href}[2]{\oldhref{#1}{\underline{#2}}}%
}

%% Patch \httplink and \httpslink to handle full URLs
% The original moderncv macros always prepend http:// or https:// to URLs.
% This causes issues when the URL already has a protocol (e.g., https://example.com)
% resulting in malformed URLs like https://https://example.com.
% This patch checks if the URL already contains "://" and uses it directly if so.
% Note: We use \str_set:Nx (x-expansion) to expand macros like \@homepage first.
\ExplSyntaxOn
\str_new:N \g_yamlresume_url_str

\RenewDocumentCommand{\httpslink}{O{}m}{
  \str_set:Nx \g_yamlresume_url_str {#2}
  \str_if_in:NnTF \g_yamlresume_url_str {://}
    { \url{#2} }
    { \url{https://#2} }
}

\RenewDocumentCommand{\httplink}{O{}m}{
  \str_set:Nx \g_yamlresume_url_str {#2}
  \str_if_in:NnTF \g_yamlresume_url_str {://}
    { \url{#2} }
    { \url{http://#2} }
}
\ExplSyntaxOff

\name{佐藤 美咲}{}
\title{マーケティングマネージャー | デジタルマーケティング戦略}
\phone[mobile]{+81 90-1234-5678}
\email{misaki.sato@example.com}
\homepage{https://www.example.com/misaki-sato}

\address{日本、東京都、渋谷区、東京都渋谷区神宮前1-2-3、150-0001}{}{}

\extrainfo{{\small \faLinkedin}\ \href{https://www.linkedin.com/in/misaki-sato}{@misaki-sato} {} {} {} • {} {} {} 
{\small \faTwitter}\ \href{https://twitter.com/misaki\_marketing}{@misaki\_marketing}}

\begin{document}

\maketitle

\section{基本情報}

\cvline{}{\begin{itemize}
\item デジタルマーケティングとブランド戦略を中心に8年以上の経験を持つマーケティングマネージャー
\item B2B/B2C両領域で、市場調査から施策立案、実行、分析までの一連のプロセスをリード
\item データドリブンな意思決定とクロスファンクショナルなチームマネジメントを得意とする
\item SEO/SEM、SNS、コンテンツマーケティング、CRMを活用した集客と顧客エンゲージメントに強み
\item マーケティング予算の最適配分とROI改善を通じて、売上成長に貢献
\end{itemize}}
\section{学歴}

\cventry{2012年4月～2016年3月}
        {学士、商学部 マーケティング専攻、成績：3.7}
        {慶應義塾大学}
        {\url{https://www.keio.ac.jp/}}
        {}
        {\begin{itemize}
\item マーケティングと経営学を中心に履修し、消費者インサイトとデータ分析の基礎を習得
\item ゼミでは企業との共同研究に参加し、新商品の市場投入戦略を提案
\item 卒業論文では「ソーシャルメディアが消費者購買行動に与える影響」をテーマに研究
\end{itemize}
\textbf{コース}：マーケティング原理、消費者行動論、市場調査、ブランド戦略、統計学、経営戦略、デジタルマーケティング}
\section{職歴}

\cventry{2021年4月～現在}
        {マーケティングマネージャー}
        {株式会社ネクストブリッジ}
        {\url{https://www.nextbridge.co.jp}}
        {}
        {\begin{itemize}
\item マーケティング部門のマネージャーとして、5名のチームをリードし年間マーケティング戦略を策定
\item 市場調査と競合分析に基づき、ターゲットセグメントの再定義とポジショニングを刷新
\item デジタル広告、SEO、コンテンツ、SNSを統合したキャンペーンを実行し、リード獲得数を前年比150\%に増加
\item マーケティング予算を四半期ごとに見直し、ROIを20\%改善
\item 営業部門と連携し、リードナーチャリング施策を構築して商談化率を向上
\end{itemize}
\textbf{キーワード}：マーケティング戦略、チームマネジメント、デジタル広告、ROI改善、リード獲得}

\cventry{2018年4月～2021年3月}
        {デジタルマーケティングスペシャリスト}
        {株式会社デジタルシフト}
        {\url{https://www.digitalshift.co.jp}}
        {}
        {\begin{itemize}
\item B2B SaaS企業において、SEO/SEM、メールマーケティング、Webinarを活用した集客を担当
\item コンテンツマーケティング戦略を立案し、オーガニックトラフィックを2年間で200\%増加
\item Google AnalyticsとCRMデータを連携し、顧客行動に基づくパーソナライズ施策を実施
\item マーケティングオートメーションを導入し、リードスコアリングとナーチャリングを自動化
\item 広告代理店と連携し、キャンペーンごとのクリエイティブ改善とA/Bテストを実施
\end{itemize}
\textbf{キーワード}：SEO、SEM、コンテンツマーケティング、MA、データ分析}

\cventry{2016年4月～2018年3月}
        {マーケティングアシスタント}
        {株式会社クリエイティブワークス}
        {\url{https://www.creativeworks.co.jp}}
        {}
        {\begin{itemize}
\item マーケティングキャンペーンの企画・運用をサポートし、SNSアカウントの運用を担当
\item プレスリリースの作成とメディアリレーションを通じて、ブランド認知度の向上に貢献
\item イベントや展示会の企画・運営を補助し、リード獲得と顧客接点の創出を支援
\item 市場調査データの収集・分析を行い、レポートを作成してチームに共有
\end{itemize}
\textbf{キーワード}：SNS運用、PR、イベント企画、市場調査}
\section{言語}

\cvline{日本語}{ネイティブ・母国語レベル \hfill \textbf{キーワード}：日本語}
\cvline{英語}{完全な実務レベル \hfill \textbf{キーワード}：TOEIC 920、ビジネス英会話}
\section{スキル}

\cvline{マーケティング戦略}{エキスパート \hfill \textbf{キーワード}：市場調査、競合分析、ブランド戦略、ポジショニング、KPI設計}
\cvline{デジタルマーケティング}{上級 \hfill \textbf{キーワード}：SEO、SEM、SNS広告、コンテンツマーケティング、メールマーケティング}
\cvline{データ分析・CRM}{上級 \hfill \textbf{キーワード}：Google Analytics、CRM、マーケティングオートメーション、A/Bテスト、データ可視化}
\cvline{マネジメント}{上級 \hfill \textbf{キーワード}：チームマネジメント、予算管理、プロジェクト管理、クロスファンクショナル連携}
\section{受賞・表彰}

\cventry{2023年12月}
        {株式会社ネクストブリッジ}
        {年間最優秀マーケティングチーム賞}
        {}
        {}
        {\begin{itemize}
\item 統合マーケティングキャンペーンの成果が評価され、社内表彰を受賞
\item リード獲得数とブランド認知度の大幅な向上に貢献
\end{itemize}}

\cventry{2017年11月}
        {日本マーケティング協会}
        {デジタルマーケティングコンテスト 優秀賞}
        {}
        {}
        {\begin{itemize}
\item 若手マーケター向けコンテストにおいて、データ活用型キャンペーンを提案し優秀賞を受賞
\end{itemize}}
\section{資格・免状}

\cventry{2023年6月}
        {Google}
        {Google Analytics 4 認定資格}
        {\url{https://skillshop.exceedlms.com/student/catalog}}
        {}
        {}

\cventry{2022年9月}
        {HubSpot Academy}
        {HubSpot コンテンツマーケティング認定}
        {\url{https://academy.hubspot.com/courses/content-marketing}}
        {}
        {}

\cventry{2015年11月}
        {日本商工会議所}
        {日商簿記検定2級}
        {}
        {}
        {}
\section{出版・著作}

\cventry{2023年3月}
        {データドリブン時代のブランド構築術}
        {マーケティングジャーナル}
        {\url{https://www.example.com/publications/data-driven-branding}}
        {}
        {\begin{itemize}
\item 顧客データと市場インサイトを活用したブランド戦略の最新手法を解説
\item B2B/B2C企業の事例を交え、実践的なフレームワークを提示
\end{itemize}}
\section{プロジェクト}

\cventry{2022年4月～2022年9月}
        {新製品の市場投入に向けた統合マーケティングキャンペーン}
        {統合マーケティングキャンペーン「Next Wave」}
        {\url{https://www.example.com/projects/next-wave}}
        {}
        {\begin{itemize}
\item ターゲット顧客の再定義からクリエイティブ制作、メディアプランニングまでを統括
\item デジタルとオフラインを組み合わせ、認知度を45\%向上、リード数を目標比130\%達成
\item キャンペーン後の顧客満足度調査でNPSが15ポイント改善
\end{itemize}
\textbf{キーワード}：統合マーケティング、新製品投入、ブランド認知}
\section{興味・関心}

\cvline{旅行}{海外旅行、温泉}
\cvline{読書}{マーケティング、ビジネス書、心理学}
\cvline{ランニング}{フルマラソン、ジョギング}
\section{ボランティア活動}

\cventry{2020年6月～現在}
        {メンター}
        {一般社団法人マーケティング支援ネットワーク}
        {\url{https://www.example.org/marketing-support}}
        {}
        {\begin{itemize}
\item スタートアップや小規模事業者向けに、デジタルマーケティングの無料相談を担当
\item マーケティング戦略の立案やSNS運用のアドバイスを通じて、地域企業の成長を支援
\item 定期的な勉強会を企画し、若手マーケターのスキルアップに貢献
\end{itemize}}
    
\end{document}